% === Begin Label Data ===
% ID: 836
% 难度星级: 
% 标签: 
% 备注: 
% 组卷引用次数: 0
% === End  Label Data ===

\begin{problem}{2021}{G}{浙江卷}{14}{解三角形}
14. (2021·浙江卷, 14)在$\triangle ABC$中，$\angle B=60^\circ$，$AB=2$，$M$是$BC$的中点，$AM=2\sqrt{3}$，则$AC=$\underline{\hspace{4em}}，$\cos\angle MAC=$\underline{\hspace{4em}}.
\end{problem}

\begin{answer}
$2\sqrt{13}$；$\dfrac{2\sqrt{39}}{13}$
\end{answer}

\begin{solutions}
方法一 ，由$\angle B=60^\circ$，$AB=2$，$AM=2\sqrt{3}$，及余弦定理可得$BM=4$，

因为$M$为$BC$的中点，所以$BC=8$. 在$\triangle ABC$中，

由余弦定理可得$AC^2=AB^2+BC^2-2BC\cdot AB\cdot\cos B=4+64-2\times 8\times 2\times\dfrac{1}{2}=52$，

所以$AC=2\sqrt{13}$，所以在$\triangle AMC$中，

由余弦定理得$\cos\angle MAC=\dfrac{AC^2+AM^2-MC^2}{2AC\cdot AM}=\dfrac{52+12-16}{2\times 2\sqrt{13}\times 2\sqrt{3}}=\dfrac{2\sqrt{39}}{13}$.

方法二，由$\angle B=60^\circ$，$AB=2$，$AM=2\sqrt{3}$，

及余弦定理可得$BM=4$，因为$M$为$BC$的中点，所以$BC=8$. 

过点$C$作$CD\perp BA$交$BA$的延长线于点$D$，则$BD=4$，$AD=2$，$CD=4\sqrt{3}$. 

所以在$\mathrm{Rt}\triangle ADC$中，$AC^2=CD^2+AD^2=48+4=52$，得$AC=2\sqrt{13}$. 

在$\triangle AMC$中，由余弦定理得$\cos\angle MAC=\dfrac{AC^2+AM^2-MC^2}{2AC\cdot AM}=\dfrac{52+12-16}{2\times 2\sqrt{13}\times 2\sqrt{3}}=\dfrac{2\sqrt{39}}{13}$.
\end{solutions}